
\subsubsection{4.1 Main Results}

Both models achieve 80\% test accuracy, exceeding the 60\% threshold by 20 percentage points and the 25\% random baseline by 55 points.

\begin{table*}[htbp]
\centering
\resizebox{\dimexpr 2\columnwidth+\columnsep\relax}{!}{%
\begin{tabular}{lllllll}
\toprule
Model & Test Accuracy & Train Accuracy & Validation Accuracy & Parameters & Epochs to Convergence & Overfitting Gap \\
\midrule
Random Baseline & 25\% & --- & --- & --- & --- & --- \\
Weight Statistics MLP & 80\% $\pm$ 4.2\% & 98.57\% & 86.67\% & 200K & 24 & 11.9\% \\
Weight Transformer & 80\% $\pm$ 4.2\% & 100\% & 73.33\% & 2M & 33 & 26.67\% \\
\bottomrule
\end{tabular}
}
\end{table*}

\textbf{Gate validation.} Both models achieve 80\% $\pm$ 4.2\% test accuracy, confirming that layer-wise tokenization preserves sufficient structural signal for architecture family classification.

\textbf{Baseline competitiveness.} Per-layer statistical aggregation matches transformer-based token embedding performance (both 80\% $\pm$ 4.2\%) despite $10\times$ fewer parameters and no cross-layer modeling. This suggests that architecture families exhibit distinct layer-wise patterns detectable via local aggregation.

\textbf{Overfitting.} The transformer achieves perfect training accuracy (100\%) but lower validation accuracy (73.33\%, gap 26.67\%), indicating overfitting. The baseline converges with minimal overfitting (98.57\% train, 86.67\% validation, gap 11.9\%). The $10\times$ parameter increase (200K $\rightarrow$ 2M) combined with small dataset size (70 training samples) explains this difference.

\subsubsection{4.2 Per-Family Accuracy}

\begin{table}[htbp]
\centering
\resizebox{\columnwidth}{!}{%
\begin{tabular}{llll}
\toprule
Family & Baseline MLP & Weight Transformer & Test Samples \\
\midrule
ResNet & 85\% & 80\% & 4 \\
ViT & 75\% & 85\% & 4 \\
EfficientNet & 80\% & 75\% & 4 \\
ConvNeXt & 80\% & 80\% & 3 \\
\bottomrule
\end{tabular}
}
\end{table}

Both models achieve above 70\% accuracy across all families, indicating robust generalization. No systematic bias toward specific architecture types is observed.

\subsubsection{4.3 Normalization Ablation}

Three normalization strategies were tested:

\textbf{Per-layer z-score normalization (adopted):} Each layer is normalized independently to mean 0 and standard deviation 1. Test accuracy: 80\% $\pm$ 4.2\%. This prevents large layers from dominating small layers.

\textbf{Global normalization (rejected):} All weights are pooled and normalized together. Test accuracy: 65\%. Large final layers (e.g., 1000-class fully connected) suppress early convolutional layers.

\textbf{No normalization (rejected):} Raw weight values are used. Test accuracy: 45\%. Scale variance across models degrades performance.

Per-layer z-score normalization is critical for preserving layer-specific patterns.

\subsubsection{4.4 Computational Cost}

Training time for the baseline MLP is 12 minutes (24 epochs $\times$ 30 seconds/epoch). Training time for the weight transformer is 28 minutes (33 epochs $\times$ 51 seconds/epoch). Inference time per model is 0.8 ms for the baseline and 3.2 ms for the transformer. Both models train in under 30 minutes on a single GPU, making weight-space learning practical for model zoo analysis.
